\documentclass[12pt,a4paper]{article}

\usepackage{cmap}
\usepackage[T2A]{fontenc}
\usepackage[utf8]{inputenc}
\usepackage[russian]{babel}
\usepackage{amsmath,amssymb}
\usepackage[a4paper,margin=2cm]{geometry}

\newcommand{\R}{\mathbb{R}}
\newcommand{\eps}{\varepsilon}
\newcommand{\dd}{\,d}
\newcommand{\mesh}{\ell}
\newcommand{\osc}{\omega}

\setlength{\parindent}{0pt}
\setlength{\parskip}{0.45em}
\sloppy

\begin{document}

\begin{center}
{\Large\bfseries Билет 6}
\end{center}

\noindent\fbox{%
  \begin{minipage}{\dimexpr\linewidth-2\fboxsep-2\fboxrule\relax}
  \normalfont
Определенный интеграл Римана. Верхние и нижние суммы Дарбу. Свойства сумм
Дарбу. Критерии интегрируемости. Интегрируемость непрерывной функции,
монотонной функции, ограниченной функции с конечным числом точек разрыва.
Аддитивность интеграла по отрезкам, линейность интеграла, интегрируемость
произведения функций, интегрируемость модуля интегрируемой функции,
интегрирование неравенств, теорема о среднем. Свойства интеграла с переменным
верхним пределом: непрерывность, дифференцируемость. Формула
Ньютона--Лейбница. Замена переменных и интегрирование по частям в определенном
интеграле.
  \end{minipage}%
}

\section*{Краткий план}

Сначала задаются разбиения отрезка, нижние и верхние суммы Дарбу и сам
интеграл Римана как общий предел этих сумм при измельчении разбиения. Затем
формулируются свойства сумм Дарбу: при измельчении нижняя сумма растет,
верхняя убывает, любая нижняя сумма не превосходит любой верхней, а разность
сумм выражается через колебания функции. Отсюда получаются критерии
интегрируемости: \(\Delta(f;T)\to0\), равенство нижнего и верхнего интегралов
Дарбу и эквивалентное определение через интегральные суммы Римана. После этого
доказываются достаточные условия интегрируемости и основные свойства
интеграла. В конце рассматривается функция
\(F(x)=\int_a^x f(t)\dd t\), формула Ньютона--Лейбница, замена переменной и
интегрирование по частям.

\section*{Определения}

\textbf{Разбиение и мелкость.}
Разбиением отрезка \([a,b]\) называется конечный набор точек
\[
  T=\{x_i\}_{i=0}^N,\qquad a=x_0<x_1<\dots<x_N=b.
\]
Его мелкость равна
\[
  \mesh(T)=\max_{1\le i\le N}(x_i-x_{i-1}).
\]
Разбиение \(T'\) называется измельчением \(T\), если \(T\subset T'\).

\textbf{Суммы Дарбу.}
Пусть \(f\) определена на \([a,b]\), а \(T=\{x_i\}_{i=0}^N\). Положим
\[
  m_i=\inf_{x\in[x_{i-1},x_i]} f(x),\qquad
  M_i=\sup_{x\in[x_{i-1},x_i]} f(x).
\]
Нижняя и верхняя суммы Дарбу:
\[
  s(f;T)=\sum_{i=1}^N (x_i-x_{i-1})m_i,\qquad
  S(f;T)=\sum_{i=1}^N (x_i-x_{i-1})M_i.
\]

\textbf{Интеграл Римана.}
Число \(J\) называется определенным интегралом Римана функции \(f\) на
\([a,b]\), если
\[
  \lim_{\mesh(T)\to0}s(f;T)=
  \lim_{\mesh(T)\to0}S(f;T)=J.
\]
Тогда пишут
\[
  J=\int_a^b f(x)\dd x,
\]
а функцию \(f\) называют интегрируемой по Риману на \([a,b]\). Необходимое
условие интегрируемости: \(f\) ограничена на \([a,b]\).

\textbf{Интегральные суммы Римана.}
Выборкой, соответствующей разбиению \(T\), называется набор
\(\xi_T=\{\xi_i\}_{i=1}^N\), где \(\xi_i\in[x_{i-1},x_i]\). Интегральная сумма
Римана:
\[
  \sigma(f;T;\xi_T)=
  \sum_{i=1}^N (x_i-x_{i-1})f(\xi_i).
\]
Для фиксированного \(T\)
\[
  s(f;T)=\inf_{\xi_T}\sigma(f;T;\xi_T),\qquad
  S(f;T)=\sup_{\xi_T}\sigma(f;T;\xi_T).
\]

\textbf{Нижний и верхний интегралы Дарбу.}
Для ограниченной функции
\[
  J_*=\sup_T s(f;T),\qquad
  J^*=\inf_T S(f;T).
\]

\textbf{Колебание.}
Колебание функции на отрезке \([x_{i-1},x_i]\):
\[
  \osc_i(f)=
  \sup_{x',x''\in[x_{i-1},x_i]}|f(x')-f(x'')|.
\]
Разность сумм Дарбу обозначается
\[
  \Delta(f;T)=S(f;T)-s(f;T).
\]

\textbf{Интеграл с переменным верхним пределом и первообразная.}
Если \(f\) интегрируема на \([a,b]\), то
\[
  F(x)=\int_a^x f(t)\dd t,\qquad x\in[a,b],
\]
называется интегралом с переменным верхним пределом. Функция \(F\) называется
первообразной функции \(f\) на \([a,b]\), если
\[
  F'(x)=f(x)\quad (x\in(a,b)),
\]
а на концах это равенство понимается через односторонние производные.

\section*{Теоремы}

\textbf{1. Свойства сумм Дарбу.}
Если \(T'\) является измельчением \(T\), то
\[
  s(f;T')\ge s(f;T),\qquad S(f;T')\le S(f;T).
\]
Для любых разбиений \(T_1,T_2\)
\[
  s(f;T_1)\le S(f;T_2).
\]
Следовательно, для ограниченной функции
\[
  s(f;T)\le J_*\le J^*\le S(f;T).
\]
Кроме того,
\[
  \Delta(f;T)=\sum_{i=1}^N (x_i-x_{i-1})\osc_i(f).
\]

\textbf{2. Критерии интегрируемости.}
Для ограниченной функции \(f\) на \([a,b]\) эквивалентны:
\[
  f\in R[a,b],
\]
\[
  \lim_{\mesh(T)\to0}\Delta(f;T)=0,
\]
\[
  J_*=J^*.
\]
Эквивалентная форма через интегральные суммы: число \(J\) равно
\(\int_a^b f(x)\dd x\) тогда и только тогда, когда
\[
  \forall\eps>0\ \exists\delta>0\ \forall T:\mesh(T)<\delta\
  \forall\xi_T\quad
  |\sigma(f;T;\xi_T)-J|<\eps.
\]

\textbf{3. Достаточные условия интегрируемости.}
На отрезке \([a,b]\) интегрируема всякая непрерывная функция, всякая монотонная
функция и всякая ограниченная функция с конечным числом точек разрыва.

\textbf{4. Основные свойства определенного интеграла.}
Если \(f\) интегрируема на \([a,b]\) и \([b,c]\), то \(f\) интегрируема на
\([a,c]\) и
\[
  \int_a^c f(x)\dd x
  =
  \int_a^b f(x)\dd x+\int_b^c f(x)\dd x.
\]
Для любого расположения точек это сохраняется при соглашениях
\[
  \int_a^a f(x)\dd x=0,\qquad
  \int_b^a f(x)\dd x=-\int_a^b f(x)\dd x.
\]
Если \(f,g\in R[a,b]\), то
\[
  \int_a^b(\alpha f+\beta g)\dd x
  =
  \alpha\int_a^b f\dd x+\beta\int_a^b g\dd x,
\]
функции \(fg\) и \(|f|\) интегрируемы,
\[
  \left|\int_a^b f(x)\dd x\right|
  \le
  \int_a^b |f(x)|\dd x,
\]
а из \(f(x)\le g(x)\) на \([a,b]\) следует
\[
  \int_a^b f(x)\dd x\le\int_a^b g(x)\dd x.
\]

\textbf{5. Теорема о среднем.}
Если \(f\) и \(g\) непрерывны на \([a,b]\), причем \(g(x)\ne0\) на
\([a,b]\), то существует \(\xi\in(a,b)\), такая что
\[
  \int_a^b f(x)g(x)\dd x
  =
  f(\xi)\int_a^b g(x)\dd x.
\]
При \(g\equiv1\) получается
\[
  \int_a^b f(x)\dd x=f(\xi)(b-a).
\]

\textbf{6. Интеграл с переменным верхним пределом.}
Если \(f\in R[a,b]\), то функция
\[
  F(x)=\int_a^x f(t)\dd t
\]
непрерывна на \([a,b]\). Если \(f\) непрерывна на \([a,b]\), то \(F\) является
первообразной \(f\):
\[
  F'(x)=f(x)\quad (x\in(a,b)),
\]
а на концах выполняются соответствующие односторонние равенства.

\textbf{7. Формула Ньютона--Лейбница.}
Если \(F\) -- первообразная непрерывной на \([a,b]\) функции \(f\), то
\[
  \int_a^b f(x)\dd x=F(b)-F(a).
\]

\textbf{8. Замена переменной и интегрирование по частям.}
Если \(\varphi\in C^1[\alpha,\beta]\), а \(f\) непрерывна на
\(\varphi([\alpha,\beta])\), то
\[
  \int_\alpha^\beta f(\varphi(t))\varphi'(t)\dd t
  =
  \int_{\varphi(\alpha)}^{\varphi(\beta)} f(x)\dd x.
\]
Если \(u,v\in C^1[a,b]\), то
\[
  \int_a^b u(x)\dd v(x)
  =
  u(x)v(x)\Big|_a^b-\int_a^b v(x)\dd u(x),
\]
то есть
\[
  \int_a^b u(x)v'(x)\dd x
  =
  u(b)v(b)-u(a)v(a)-\int_a^b v(x)u'(x)\dd x.
\]

\section*{Доказательства}

\textbf{Свойства сумм Дарбу.}
Пусть \(T'\) получается из \(T\) добавлением одной точки \(x'\in(x_{j-1},x_j)\).
На меньших отрезках инфимум не меньше инфимума на \([x_{j-1},x_j]\), поэтому
вклад в нижнюю сумму не уменьшается:
\[
  s(f;T')-s(f;T)\ge0.
\]
Для супремумов наоборот: на меньших отрезках супремум не больше супремума на
исходном отрезке, значит
\[
  S(f;T')-S(f;T)\le0.
\]
Если добавлено несколько точек, применяем этот шаг последовательно.

Для двух произвольных разбиений берем их объединение \(T=T_1\cup T_2\). Тогда
по уже доказанному
\[
  s(f;T_1)\le s(f;T)\le S(f;T)\le S(f;T_2).
\]
Отсюда следуют неравенства
\[
  s(f;T)\le J_*\le J^*\le S(f;T).
\]
Наконец, так как
\[
  \osc_i(f)=M_i-m_i,
\]
то
\[
  \Delta(f;T)
  =
  \sum_{i=1}^N(x_i-x_{i-1})(M_i-m_i)
  =
  \sum_{i=1}^N(x_i-x_{i-1})\osc_i(f).
\]

\textbf{Критерий через \(\Delta(f;T)\).}
Если \(f\) интегрируема и
\[
  \int_a^b f(x)\dd x=J,
\]
то при достаточно малой мелкости
\[
  |s(f;T)-J|<\frac{\eps}{2},\qquad
  |S(f;T)-J|<\frac{\eps}{2}.
\]
Следовательно,
\[
  \Delta(f;T)=S(f;T)-s(f;T)
  \le |S(f;T)-J|+|s(f;T)-J|<\eps.
\]

Обратно, пусть \(\Delta(f;T)\to0\). Тогда \(f\) ограничена: иначе одна из сумм
Дарбу была бы бесконечной и \(\Delta(f;T)\) не могла бы стремиться к нулю.
По свойству
\[
  s(f;T)\le J_*\le J^*\le S(f;T)
\]
получаем
\[
  0\le J^*-J_*\le S(f;T)-s(f;T)=\Delta(f;T).
\]
Значит \(J_*=J^*=:J\). Тогда
\[
  s(f;T)\le J\le S(f;T),
\]
и из \(\Delta(f;T)\to0\) следует
\[
  s(f;T)\to J,\qquad S(f;T)\to J.
\]
По определению \(f\) интегрируема и \(\int_a^b f=J\).

\textbf{Критерий через интегральные суммы.}
Для фиксированного разбиения \(T\)
\[
  s(f;T)\le \sigma(f;T;\xi_T)\le S(f;T)
\]
при любой выборке. Поэтому сходимость нижней и верхней сумм Дарбу к \(J\)
равносильна равномерной по выборкам сходимости
\[
  \sigma(f;T;\xi_T)\to J\qquad(\mesh(T)\to0).
\]
Точно так же из сходимости всех интегральных сумм к \(J\), беря инфимум и
супремум по выборкам, получаем сходимость \(s(f;T)\) и \(S(f;T)\) к \(J\).

\textbf{Интегрируемость непрерывной функции.}
Если \(f\) непрерывна на \([a,b]\), то по теореме Кантора она равномерно
непрерывна. Пусть
\[
  \omega(\rho)=
  \sup_{\substack{x,x'\in[a,b]\\ |x-x'|\le\rho}}|f(x)-f(x')|.
\]
Тогда \(\omega(\rho)\to0\) при \(\rho\to0\). Для любого разбиения
\[
  \osc_i(f)\le\omega(\mesh(T)),
\]
поэтому
\[
  \Delta(f;T)
  =
  \sum_{i=1}^N(x_i-x_{i-1})\osc_i(f)
  \le
  (b-a)\omega(\mesh(T))\to0.
\]
По критерию Дарбу \(f\) интегрируема.

\textbf{Интегрируемость монотонной функции.}
Пусть \(f\) не убывает на \([a,b]\). Тогда
\[
  \osc_i(f)=f(x_i)-f(x_{i-1}),
\]
и
\[
  \Delta(f;T)
  =
  \sum_{i=1}^N(x_i-x_{i-1})(f(x_i)-f(x_{i-1}))
  \le
  \mesh(T)\sum_{i=1}^N(f(x_i)-f(x_{i-1})).
\]
Последняя сумма телескопируется:
\[
  \Delta(f;T)\le \mesh(T)(f(b)-f(a))\to0.
\]
Случай невозрастающей функции сводится к этому применением рассуждения к
\(-f\).

\textbf{Ограниченная функция с конечным числом разрывов.}
Пусть \(|f|\le M\), а точки разрыва образуют конечное множество
\(\{c_1,\dots,c_k\}\). Для заданного \(\eps>0\) покроем эти точки конечным
набором интервалов суммарной длины меньше \(\eps/(4M)\). На оставшейся части
отрезка функция непрерывна на конечном объединении замкнутых отрезков, значит
равномерно непрерывна; выбираем \(\delta\) так, чтобы вне выбранных малых
интервалов колебание на отрезках длины \(<\delta\) было меньше
\(\eps/(2(b-a))\).

Если \(\mesh(T)<\delta\), то вклад отрезков разбиения, попавших в малые
окрестности точек разрыва, не превосходит \(2M\) умножить на их суммарную
длину, то есть меньше \(\eps/2\). На остальных отрезках суммарный вклад меньше
\(\eps/2\). Следовательно,
\[
  \Delta(f;T)<\eps,
\]
и функция интегрируема по критерию Дарбу.

\textbf{Линейность.}
Для любых разбиения и выборки
\[
  \sigma(\alpha f+\beta g;T;\xi_T)
  =
  \alpha\sigma(f;T;\xi_T)+\beta\sigma(g;T;\xi_T).
\]
Переходя к пределу при \(\mesh(T)\to0\), получаем интегрируемость
\(\alpha f+\beta g\) и формулу линейности интеграла.

\textbf{Интегрирование неравенств.}
Если \(f(x)\le g(x)\) на \([a,b]\), то для любых \(T,\xi_T\)
\[
  \sigma(f;T;\xi_T)\le\sigma(g;T;\xi_T).
\]
Переход к пределу по интегральным суммам дает
\[
  \int_a^b f(x)\dd x\le\int_a^b g(x)\dd x.
\]

\textbf{Интегрируемость модуля.}
Из неравенства
\[
  \bigl||f(x')|-|f(x'')|\bigr|\le |f(x')-f(x'')|
\]
следует
\[
  \osc_i(|f|)\le\osc_i(f).
\]
Поэтому
\[
  0\le\Delta(|f|;T)\le\Delta(f;T)\to0.
\]
Значит \(|f|\) интегрируема. Так как
\[
  -|f(x)|\le f(x)\le |f(x)|,
\]
из интегрирования неравенств и линейности следует
\[
  \left|\int_a^b f(x)\dd x\right|
  \le
  \int_a^b |f(x)|\dd x.
\]

\textbf{Интегрируемость произведения.}
Интегрируемые функции ограничены; пусть
\[
  M_f=\sup_{[a,b]}|f|,\qquad M_g=\sup_{[a,b]}|g|.
\]
Тогда
\[
  |f(x)g(x)-f(x')g(x')|
  \le
  M_g|f(x)-f(x')|+M_f|g(x)-g(x')|.
\]
Следовательно,
\[
  \osc_i(fg)\le M_g\osc_i(f)+M_f\osc_i(g),
\]
и
\[
  \Delta(fg;T)
  \le
  M_g\Delta(f;T)+M_f\Delta(g;T)\to0.
\]
По критерию Дарбу произведение \(fg\) интегрируемо.

\textbf{Аддитивность по отрезкам.}
Пусть \(f\) интегрируема на \([a,b]\) и \([b,c]\). Возьмем разбиение
\([a,c]\) и добавим к нему точку \(b\). При измельчении разбиения
интегральная сумма на \([a,c]\) распадается в сумму интегральных сумм на
\([a,b]\) и \([b,c]\), а добавление одной точки не меняет предела при
\(\mesh(T)\to0\). Поэтому
\[
  \int_a^c f(x)\dd x
  =
  \int_a^b f(x)\dd x+\int_b^c f(x)\dd x.
\]
Формулы для произвольного порядка точек следуют из соглашений
\(\int_a^a f=0\) и \(\int_b^a f=-\int_a^b f\).

\textbf{Теорема о среднем.}
Пусть \(f,g\) непрерывны и \(g\ne0\) на \([a,b]\). Тогда \(g\) сохраняет знак,
поэтому
\[
  \int_a^b g(x)\dd x\ne0.
\]
Положим
\[
  \Phi(x)=\int_a^x f(t)g(t)\dd t,\qquad
  G(x)=\int_a^x g(t)\dd t.
\]
По свойству интеграла с переменным верхним пределом
\[
  \Phi'(x)=f(x)g(x),\qquad G'(x)=g(x).
\]
По теореме Коши о среднем для функций одной переменной существует
\(\xi\in(a,b)\), такая что
\[
  \frac{\Phi(b)-\Phi(a)}{G(b)-G(a)}
  =
  \frac{\Phi'(\xi)}{G'(\xi)}
  =
  f(\xi).
\]
Значит
\[
  \int_a^b f(x)g(x)\dd x
  =
  f(\xi)\int_a^b g(x)\dd x.
\]

\textbf{Интеграл с переменным верхним пределом.}
Если \(f\in R[a,b]\), то \(f\) ограничена: \(|f|\le M\). Для
\[
  F(x)=\int_a^x f(t)\dd t
\]
и \(x_1,x_2\in[a,b]\) по аддитивности
\[
  F(x_2)-F(x_1)=\int_{x_1}^{x_2}f(t)\dd t.
\]
Следовательно,
\[
  |F(x_2)-F(x_1)|
  \le
  \int_{x_1}^{x_2}|f(t)|\dd t
  \le
  M|x_2-x_1|.
\]
Значит \(F\) непрерывна.

Если \(f\) непрерывна, то при \(x\ne x_0\)
\[
  \frac{F(x)-F(x_0)}{x-x_0}-f(x_0)
  =
  \frac{1}{x-x_0}\int_{x_0}^x\bigl(f(t)-f(x_0)\bigr)\dd t.
\]
По непрерывности \(f\) в точке \(x_0\) правая часть по модулю не превосходит
\(\eps\) при достаточно малом \(|x-x_0|\). Поэтому
\[
  F'(x_0)=f(x_0)
\]
во внутренних точках, а на концах -- в одностороннем смысле.

\textbf{Формула Ньютона--Лейбница.}
Пусть \(F\) -- произвольная первообразная непрерывной функции \(f\). Функция
\[
  \Phi(x)=\int_a^x f(t)\dd t
\]
тоже является первообразной \(f\). Две первообразные на отрезке отличаются на
константу, поэтому
\[
  F(x)=\Phi(x)+C.
\]
Тогда
\[
  F(b)-F(a)=\Phi(b)-\Phi(a)=\int_a^b f(x)\dd x.
\]

\textbf{Замена переменной.}
Пусть \(A\) -- первообразная функции \(f\) на \(\varphi([\alpha,\beta])\).
Тогда по правилу дифференцирования сложной функции
\[
  \frac{\dd}{\dd t}A(\varphi(t))
  =
  f(\varphi(t))\varphi'(t).
\]
Применяя формулу Ньютона--Лейбница к левой части, получаем
\[
  \int_\alpha^\beta f(\varphi(t))\varphi'(t)\dd t
  =
  A(\varphi(\beta))-A(\varphi(\alpha))
  =
  \int_{\varphi(\alpha)}^{\varphi(\beta)} f(x)\dd x.
\]

\textbf{Интегрирование по частям.}
Так как
\[
  (uv)'=u'v+uv',
\]
то по формуле Ньютона--Лейбница и линейности
\[
  u(b)v(b)-u(a)v(a)
  =
  \int_a^b (uv)'(x)\dd x
  =
  \int_a^b u'(x)v(x)\dd x+\int_a^b u(x)v'(x)\dd x.
\]
Перенося первый интеграл вправо, получаем
\[
  \int_a^b u(x)v'(x)\dd x
  =
  u(x)v(x)\Big|_a^b-\int_a^b v(x)u'(x)\dd x.
\]

\section*{Финальная устная версия}

Разбиение отрезка \([a,b]\) -- это
\[
  a=x_0<x_1<\dots<x_N=b,\qquad
  \mesh(T)=\max_i(x_i-x_{i-1}).
\]
Для функции \(f\) на каждом отрезке разбиения берут
\[
  m_i=\inf_{[x_{i-1},x_i]}f,\qquad
  M_i=\sup_{[x_{i-1},x_i]}f
\]
и определяют суммы Дарбу
\[
  s(f;T)=\sum_i (x_i-x_{i-1})m_i,\qquad
  S(f;T)=\sum_i (x_i-x_{i-1})M_i.
\]
Число \(J\) называется интегралом Римана, если при \(\mesh(T)\to0\)
\[
  s(f;T)\to J,\qquad S(f;T)\to J.
\]
Тогда \(J=\int_a^b f(x)\dd x\), а функция интегрируема.

Главные свойства сумм Дарбу: при измельчении разбиения нижняя сумма не
уменьшается, а верхняя не увеличивается:
\[
  s(f;T')\ge s(f;T),\qquad S(f;T')\le S(f;T).
\]
Для любых \(T_1,T_2\) имеем \(s(f;T_1)\le S(f;T_2)\). Если
\[
  J_*=\sup_T s(f;T),\qquad J^*=\inf_T S(f;T),
\]
то
\[
  s(f;T)\le J_*\le J^*\le S(f;T).
\]
Разность сумм выражается через колебания:
\[
  \Delta(f;T)=S(f;T)-s(f;T)
  =
  \sum_i (x_i-x_{i-1})\osc_i(f).
\]
Критерий интегрируемости:
\[
  f\in R[a,b]\quad\Longleftrightarrow\quad
  \Delta(f;T)\to0\quad(\mesh(T)\to0),
\]
что равносильно \(J_*=J^*\). То же самое можно сказать через интегральные
суммы
\[
  \sigma(f;T;\xi_T)=\sum_i (x_i-x_{i-1})f(\xi_i):
\]
они должны стремиться к одному и тому же числу при \(\mesh(T)\to0\) равномерно
по выборкам.

Непрерывная функция интегрируема, потому что на отрезке она равномерно
непрерывна и колебания на мелких отрезках малы:
\[
  \Delta(f;T)\le (b-a)\omega(\mesh(T))\to0.
\]
Монотонная функция интегрируема, потому что для неубывающей функции
\[
  \Delta(f;T)\le \mesh(T)(f(b)-f(a))\to0.
\]
Ограниченная функция с конечным числом разрывов интегрируема: малые
окрестности разрывов дают малый вклад из-за ограниченности, а вне них функция
равномерно непрерывна на конечном объединении отрезков.

Интеграл аддитивен по отрезкам:
\[
  \int_a^c f=\int_a^b f+\int_b^c f.
\]
Он линеен:
\[
  \int_a^b(\alpha f+\beta g)=
  \alpha\int_a^b f+\beta\int_a^b g.
\]
Если \(f,g\) интегрируемы, то \(fg\) интегрируема, потому что
\[
  \osc_i(fg)\le M_g\osc_i(f)+M_f\osc_i(g).
\]
Также интегрируема \(|f|\), так как
\[
  \osc_i(|f|)\le\osc_i(f),
\]
и
\[
  \left|\int_a^b f\right|\le \int_a^b |f|.
\]
Если \(f\le g\), то
\[
  \int_a^b f\le\int_a^b g.
\]
Теорема о среднем: если \(f,g\) непрерывны и \(g\ne0\) на \([a,b]\), то для
некоторого \(\xi\in(a,b)\)
\[
  \int_a^b f(x)g(x)\dd x=f(\xi)\int_a^b g(x)\dd x.
\]

Для
\[
  F(x)=\int_a^x f(t)\dd t
\]
если \(f\) интегрируема, то \(F\) непрерывна, потому что
\[
  |F(x_2)-F(x_1)|\le M|x_2-x_1|.
\]
Если \(f\) непрерывна, то \(F'(x)=f(x)\), поскольку
\[
  \frac{F(x)-F(x_0)}{x-x_0}-f(x_0)
  =
  \frac{1}{x-x_0}\int_{x_0}^x(f(t)-f(x_0))\dd t\to0.
\]
Отсюда следует формула Ньютона--Лейбница:
\[
  \int_a^b f(x)\dd x=F(b)-F(a)
\]
для любой первообразной \(F\) непрерывной функции \(f\).

Если \(\varphi\in C^1[\alpha,\beta]\), то
\[
  \int_\alpha^\beta f(\varphi(t))\varphi'(t)\dd t
  =
  \int_{\varphi(\alpha)}^{\varphi(\beta)} f(x)\dd x,
\]
потому что первообразная \(A\) функции \(f\) дает
\((A\circ\varphi)'=f(\varphi)\varphi'\). Если \(u,v\in C^1[a,b]\), то из
\((uv)'=u'v+uv'\) и формулы Ньютона--Лейбница получается интегрирование по
частям:
\[
  \int_a^b u(x)v'(x)\dd x
  =
  u(x)v(x)\Big|_a^b-\int_a^b v(x)u'(x)\dd x.
\]

\section*{Источники}

Иванов Г.Е. \emph{Лекции по математическому анализу. Часть 1}:
\texttt{IvGE\_1.pdf}, стр. 207--227.

Основные роли источника: стр. 207--213 -- суммы Дарбу, определение интеграла
Римана, свойства сумм Дарбу и критерий интегрируемости; стр. 213--215 --
интегральные суммы Римана и эквивалентное определение интеграла; стр. 215--223
-- линейность, интегрирование неравенств, интегрируемость модуля, изменение
значений в конечном числе точек, аддитивность по отрезкам, интегрируемость
непрерывных, кусочно-непрерывных, монотонных функций и произведения; стр.
223--227 -- интеграл с переменным верхним пределом, формула
Ньютона--Лейбница, замена переменной, интегрирование по частям и теорема о
среднем.

Дополнительно был просмотрен \texttt{IvGE\_2.pdf}: прямого материала по
интегралу Римана одной переменной для этого билета там нет; найденные
совпадения относятся к более поздним темам.

\end{document}
